% Vista científica exacta materializada desde una rama condicional.
% Fuente: source/demostraciones_primarias/09_06h_atlas_familias_rev6.tex; rama: HMTIncludeAtlasPhysical.
% No modifica el contenido matemático de la rama de origen.
\chapter[Correspondencia con sectores físicos]{Correspondencia entre clases internas y sectores físicos}
\label{sec:censo-fisico-atlas-rev7}

La realización física conserva los cocientes anteriores y añade los datos
que una representación necesita sin convertirlos en nuevas celdas. Para una
sección persistente \(X\), escribimos de forma tipada
\[
 \mathcal P_X=
 \bigl[q(X),\,[X]_{\mathcal R},\,\operatorname{Rep}(X),\,
       \gamma_X,\,\tau_{12}(X),\,\sigma_X\bigr].
\]
Aquí \(q(X)\) fija la multisección, \([X]_{\mathcal R}\) la familia de ruta,
\(\operatorname{Rep}(X)\) la representación, \(\gamma_X\) la historia
persistente, \(\tau_{12}(X)\) su palabra dodecafásica y \(\sigma_X\) la firma
entera extraída después del registro. Los niveles \(G0,\ldots,G4\) siguen
siendo profundidades del atlas; no se renombran como las tres generaciones
físicas.

La información de sabor se conserva mediante tres coordenadas distintas:
\[
 \operatorname{flav}(X)=
 \bigl(\sigma_{ud}(X),g(X),\chi_{\rm fina}(X)\bigr).
\]
La primera separa las ramas leptónica, neutrino, tipo arriba y tipo abajo;
la segunda distingue la generación física; la tercera conserva la
orientación y la memoria que separan secciones dentro de una misma rama. En
los quarks, la órbita residual \(R/G/B\) se añade después de fijar el sabor y
antes de evaluar el carácter. Color, sabor y masa son, por tanto, lecturas de
tipos diferentes.

% El inventario humano sectorial se conserva en su totalidad,
% pero se monta después del cierre electrónico, en la residencia sectorial.

Este censo no agota las realizaciones de la ley. Mesones y bariones se
construyen componiendo registros enriquecidos antes de proyectar la firma;
resonancias y partículas virtuales se distinguen por su horizonte de
persistencia. Los sectores Fibonacci, \(Z_6\) y Majorana--Ising conservan
codominios topológicos propios y no se fuerzan dentro de una tabla ordinaria
de masas.

\section{Firmas de trayectoria y pares dodecafásicos coindexados por sabor}

Para el conjunto de sabores
\(\mathcal Q=\{u,d,c,s,t,b\}\), la tabla define una correspondencia finita
entre una huella geométrica y un par dodecafásico. En cada fila se cumple
\[
 W_\pm=6n_A\pm s_C.
\]
Esta tabla no es una comparación a posteriori, sino la gráfica de un mapa
finito sobre el dominio nominal certificado. Si
\(\mathscr R_{\mathcal Q}^{\rm nom}=\{\rho_u,\rho_d,\rho_c,\rho_s,\rho_t,\rho_b\}\)
es el conjunto de huellas CT--108 que siguen, definimos
\[
 H_{\rm flav}:\mathscr R_{\mathcal Q}^{\rm nom}\longrightarrow\mathbb Z^2,
 \qquad H_{\rm flav}(\rho_q)=(n_A(q),s_C(q)).
\]
Las seis imágenes son únicas, la tabla exhausta el dominio y ninguna columna
contiene una masa. Por ello \(H_{\rm flav}\) es un extractor causal cerrado en
el dominio nominal: primero lee huella, orientación y memoria, y sólo después
el carácter evalúa la masa. La órbita ternaria de color se cocienta sin cambiar
\(H_{\rm flav}\); la reversión conserva la clase de masa y transporta la
representación partícula--antipartícula. La tabla es, por tanto, la definición
ejecutable del morfismo y no una obligación futura:

\begin{longtable}{c c c c c c c}
\toprule
Sabor & Generación & Giros/cambios & Eje & \(N,E,S,O\)
& \((n_A,s_C)\) & \((W_+,W_-)\)\\
\midrule
\endfirsthead
\toprule
Sabor & Generación & Giros/cambios & Eje & \(N,E,S,O\)
& \((n_A,s_C)\) & \((W_+,W_-)\)\\
\midrule
\endhead
\(u\)&1&21/30&N--S&34,31,31,12&(11,7)&(73,59)\\
\(d\)&1&20/32&E--O&27,35,18,28&(17,8)&(110,94)\\
\(c\)&2&21/29&E--O&20,38,24,26&(61,8)&(374,358)\\
\(s\)&2&17/27&N--S&50,32,12,14&(41,-3)&(243,249)\\
\(t\)&3&21/26&N--S&46,27,16,19&(100,-1)&(599,601)\\
\(b\)&3&21/30&E--O&12,63,9,24&(71,-5)&(421,431)\\
\bottomrule
\end{longtable}

La coincidencia de dos giros o de una altura espectral no identifica
especies: la ruta cardinal, la orientación, el registro dodecafásico y la
representación permanecen en el objeto. La cadena
celda--ruta--registro cronológico enriquecido--firma--carácter se define con
independencia de \(H_{\rm flav}\). Las
cuatro orientaciones iniciales de cada celda producen
\(81\cdot4=324\) genealogías ejecutables; ese cardinal cuenta historias, no
partículas, familias ni masas.
